\documentclass[12pt,a4paper]{article}

% Packages
\usepackage[utf8]{inputenc}
\usepackage[english]{babel}
\usepackage{amsmath}
\usepackage{amsfonts}
\usepackage{amssymb}
\usepackage{graphicx}
\usepackage{hyperref}
\usepackage{cite}

% Title and author information
\title{An Example Article: Introduction to LaTeX Document Preparation}
\author{John Doe\\
        Department of Computer Science\\
        Example University\\
        \texttt{john.doe@example.edu}}
\date{\today}

\begin{document}

\maketitle

\begin{abstract}

Changing something more than this

This article serves as an example of a well-structured LaTeX document. It demonstrates the use of various LaTeX features including document structure, mathematical equations, figures, tables, and references. The purpose of this document is to provide a template for academic writing and to showcase the capabilities of the LaTeX typesetting system. Through various sections, we explore fundamental concepts and best practices in document preparation using LaTeX.
\end{abstract}

\section{Introduction}
\label{sec:introduction}

LaTeX is a high-quality typesetting system that is widely used for technical and scientific documents~\cite{lamport1994}. It is built on top of TeX, the powerful typesetting engine created by Donald Knuth~\cite{knuth1984}. LaTeX provides excellent support for mathematical notation, cross-referencing, and bibliography management. Unlike word processors, LaTeX separates content from formatting, allowing authors to focus on the structure and meaning of their text rather than its appearance.

The popularity of LaTeX in academic circles stems from several key advantages:
\begin{itemize}
    \item Professional typesetting quality
    \item Excellent mathematical formula support
    \item Automated reference and citation management
    \item Consistent formatting throughout the document
    \item Platform independence and long-term stability
\end{itemize}

This article is organized as follows: Section~\ref{sec:methodology} describes the methodology, Section~\ref{sec:results} presents the results, and Section~\ref{sec:conclusion} concludes the paper.

\section{Methodology}
\label{sec:methodology}

In this section, we demonstrate various LaTeX features that are commonly used in academic writing. We begin by showing how to include mathematical equations and formulas.

\subsection{Mathematical Equations}
\label{subsec:equations}

LaTeX excels at typesetting mathematical content. For example, inline mathematics can be written as $E = mc^2$, which is Einstein's famous mass-energy equivalence formula.

More complex equations can be displayed on separate lines:
\begin{equation}
    \int_{-\infty}^{\infty} e^{-x^2} dx = \sqrt{\pi}
    \label{eq:gaussian}
\end{equation}

The Gaussian integral in Equation~\ref{eq:gaussian} is a fundamental result in calculus. We can also write systems of equations:
\begin{align}
    \nabla \cdot \mathbf{E} &= \frac{\rho}{\epsilon_0} \label{eq:maxwell1}\\
    \nabla \cdot \mathbf{B} &= 0 \label{eq:maxwell2}\\
    \nabla \times \mathbf{E} &= -\frac{\partial \mathbf{B}}{\partial t} \label{eq:maxwell3}\\
    \nabla \times \mathbf{B} &= \mu_0\mathbf{J} + \mu_0\epsilon_0\frac{\partial \mathbf{E}}{\partial t} \label{eq:maxwell4}
\end{align}

Equations~\ref{eq:maxwell1}--\ref{eq:maxwell4} are Maxwell's equations, which form the foundation of classical electromagnetism.

\subsection{Tables and Figures}
\label{subsec:tables}

Tables are essential for presenting data in a structured format. Table~\ref{tab:example} shows an example of a simple table.

\begin{table}[h]
\centering
\begin{tabular}{|l|c|r|}
\hline
\textbf{Item} & \textbf{Quantity} & \textbf{Price (\$)} \\
\hline
Widget A & 10 & 25.00 \\
Widget B & 5 & 42.50 \\
Widget C & 15 & 18.75 \\
\hline
\textbf{Total} & \textbf{30} & \textbf{86.25} \\
\hline
\end{tabular}
\caption{Example table showing product inventory and pricing.}
\label{tab:example}
\end{table}

\section{Results}
\label{sec:results}

The results of our investigation demonstrate the versatility and power of LaTeX for document preparation. We found that LaTeX provides superior control over document formatting compared to traditional word processors, particularly for documents containing complex mathematical notation.

\subsection{Key Findings}

Our analysis revealed several important findings:
\begin{enumerate}
    \item LaTeX produces consistently formatted documents across different platforms
    \item The learning curve is steep initially but becomes manageable with practice
    \item Version control systems work seamlessly with LaTeX source files
    \item Large documents can be efficiently managed through modular structure
\end{enumerate}

\subsection{Discussion}

The advantages of using LaTeX become particularly evident in collaborative academic writing. The plain text format enables easy version control using systems like Git, and the separation of content from style ensures consistency across contributions from multiple authors.

Furthermore, the extensive package ecosystem provides solutions for virtually any typesetting requirement, from complex diagrams to specialized notation in various scientific fields. For comprehensive guidance on advanced LaTeX features and packages, readers are referred to The LaTeX Companion~\cite{mittelbach2004}.

\section{Conclusion}
\label{sec:conclusion}

This article has demonstrated the basic structure and capabilities of LaTeX for academic document preparation. We have shown how to create structured documents with sections, subsections, mathematical equations, tables, and references. LaTeX remains an essential tool for researchers and academics who require high-quality typesetting and precise control over their document formatting.

For those beginning their journey with LaTeX, we recommend starting with simple documents and gradually incorporating more advanced features. The investment in learning LaTeX pays dividends in the form of professional-quality documents and efficient workflow management.

\subsection{Future Work}

Future enhancements to this example could include:
\begin{itemize}
    \item Adding actual figures and graphics using the \texttt{graphicx} package
    \item Using BibTeX for automated bibliography management
    \item Demonstrating advanced table formatting with packages like \texttt{booktabs}
    \item Showing code listings with syntax highlighting
    \item Including cross-references to other documents
\end{itemize}

\section*{Acknowledgments}

The author would like to thank the LaTeX community for their continuous efforts in maintaining and improving this excellent typesetting system.

\begin{thebibliography}{9}
\bibitem{lamport1994}
Leslie Lamport. 
\textit{LaTeX: A Document Preparation System}. 
Addison-Wesley, 2nd edition, 1994.

\bibitem{knuth1984}
Donald E. Knuth. 
\textit{The TeXbook}. 
Addison-Wesley, 1984.

\bibitem{mittelbach2004}
Frank Mittelbach and Michel Goossens. 
\textit{The LaTeX Companion}. 
Addison-Wesley, 2nd edition, 2004.
\end{thebibliography}

\end{document}
